\documentclass{article}
\usepackage{amsmath,amssymb,amsthm}
\usepackage{algorithm}
\usepackage[noend]{algpseudocode}
\usepackage{enumitem}
\usepackage{hyperref}
\usepackage{geometry}
\geometry{margin=1in}
\newtheorem{assumption}{Assumption}
\newtheorem{theorem}{Theorem}
\newtheorem{lemma}{Lemma}
\newcommand{\EE}{\mathbb{E}}
\newcommand{\RR}{\mathbb{R}}

\begin{document}

\section*{Algorithm Name}
\textbf{Stochastic Gradient Descent with Warm Restarts (SGDR).}

\section*{Mathematical Formulation}
Let $f:\RR^{d}\rightarrow\RR$ be the (possibly non‑convex) objective function and let
$\{z_{t}\}_{t\ge 0}$ be an i.i.d.\ data stream.  
At iteration $t$ we compute a stochastic gradient
\[
g_{t}\;:=\;\nabla f(w_{t};z_{t})\in\RR^{d}.
\]
The learning‑rate schedule is a cosine annealing schedule that is reset every $M$ steps:
\[
\eta_{t}\;=\;\eta_{\max}\,\frac{1}{2}\Bigl(1+\cos\!\bigl(\pi\,\frac{t-s_{k}}{M}\bigr)\Bigr),
\qquad 
s_{k}=kM\;(k=0,1,2,\dots)
\tag{1}
\]
where $s_{k}$ is the start index of the $k$‑th cycle.  
The update rule is
\[
w_{t+1}\;=\;w_{t}-\eta_{t}\,g_{t},\qquad t=0,1,2,\dots
\tag{2}
\]

\section*{Pseudocode}
\begin{algorithm}[H]
\caption{SGDR}
\begin{algorithmic}[1]
\State \textbf{Input:} initial point $w_{0}\in\RR^{d}$, maximal step size $\eta_{\max}>0$, cycle length $M\in\mathbb{N}$, total iterations $T$.
\State Set $t\gets0$.
\While{$t<T$}
    \State $k\gets \lfloor t/M\rfloor$ \Comment{current cycle index}
    \State $s_{k}\gets kM$ \Comment{start of current cycle}
    \State $\displaystyle \eta_{t}\gets \eta_{\max}\frac12\bigl(1+\cos(\pi\frac{t-s_{k}}{M})\bigr)$ \Comment{cosine annealing}
    \State Sample $z_{t}$ and compute stochastic gradient $g_{t}\gets\nabla f(w_{t};z_{t})$.
    \State $w_{t+1}\gets w_{t}-\eta_{t}\,g_{t}$.
    \State $t\gets t+1$.
\EndWhile
\State \textbf{Output:} $\{w_{t}\}_{t=0}^{T}$ (or the averaged iterate $\bar w_{T}=\frac1T\sum_{t=0}^{T-1}w_{t}$).
\end{algorithmic}
\end{algorithm}

\section*{Dry Run Example}
Consider the scalar quadratic $f(w)=(w-3)^{2}$, thus $\nabla f(w)=2(w-3)$.  
Set $w_{0}=0$, $\eta_{\max}=0.1$, $M=4$, and perform three iterations ($T=3$).  
Because $t<M$, the cosine term equals $\frac12(1+\cos(\pi t/M))$, yielding
\[
\eta_{0}=0.1,\;\eta_{1}=0.1\frac12\!\bigl(1+\cos(\tfrac{\pi}{4})\bigr)\approx0.0924,\;
\eta_{2}=0.1\frac12\!\bigl(1+\cos(\tfrac{\pi}{2})\bigr)=0.05.
\]

\begin{enumerate}[label=\textbf{Iter \arabic*:}, leftmargin=*]
\item $t=0$:
\[
g_{0}=2(w_{0}-3)=2(0-3)=-6,\qquad 
w_{1}=w_{0}-\eta_{0}g_{0}=0-0.1(-6)=0.6,
\]
\[
f(w_{1})=(0.6-3)^{2}=5.76.
\]

\item $t=1$:
\[
g_{1}=2(w_{1}-3)=2(0.6-3)=-4.8,\qquad 
w_{2}=w_{1}-\eta_{1}g_{1}=0.6-0.0924(-4.8)\approx1.045.
\]
\[
f(w_{2})=(1.045-3)^{2}\approx3.80.
\]

\item $t=2$:
\[
g_{2}=2(w_{2}-3)=2(1.045-3)=-3.91,\qquad 
w_{3}=w_{2}-\eta_{2}g_{2}=1.045-0.05(-3.91)\approx1.241.
\]
\[
f(w_{3})=(1.241-3)^{2}\approx3.09.
\]
\end{enumerate}
The objective value decreases monotonically, illustrating the effect of the decaying step size.

\newpage
\section*{Assumptions}
\begin{assumption}[Smoothness]\label{as:smooth}
$f$ is $L$‑smooth, i.e.
\[
\|\nabla f(x)-\nabla f(y)\|\le L\|x-y\|\quad\forall x,y\in\RR^{d}.
\]
Equivalently,
\[
f(y)\le f(x)+\langle\nabla f(x),y-x\rangle+\frac{L}{2}\|y-x\|^{2}.
\tag{3}
\]
\end{assumption}

\begin{assumption}[Unbiased Stochastic Gradient]\label{as:unbiased}
For every $t$, 
\[
\EE\!\bigl[g_{t}\mid w_{t}\bigr]=\nabla f(w_{t}).
\]
\end{assumption}

\begin{assumption}[Bounded Variance]\label{as:variance}
There exists $\sigma^{2}\ge0$ such that
\[
\EE\!\bigl[\|g_{t}-\nabla f(w_{t})\|^{2}\mid w_{t}\bigr]\le\sigma^{2}\quad\forall t.
\tag{4}
\]
\end{assumption}

\begin{assumption}[Step‑size Bounds]\label{as:stepsize}
The cosine schedule \eqref{1} satisfies
\[
0<\eta_{\min}\le\eta_{t}\le\eta_{\max}\le\frac{1}{L},
\qquad\forall t,
\tag{5}
\]
where $\eta_{\min}=\eta_{\max}\frac12\bigl(1+\cos(\pi)\bigr)=0$ is allowed, but we will use the upper bound $\eta_{\max}\le1/L$ for stability.
\end{assumption}

\section*{Main Convergence Theorem}
\begin{theorem}\label{thm:main}
Assume \ref{as:smooth}–\ref{as:stepsize}. Let $w_{0}$ be arbitrary and run SGDR for $T$ iterations. Define the averaged iterate
\[
\bar w_{T}\;:=\;\frac{1}{T}\sum_{t=0}^{T-1}w_{t}.
\]
Then
\[
\frac{1}{T}\sum_{t=0}^{T-1}\EE\!\bigl[\|\nabla f(w_{t})\|^{2}\bigr]
\;\le\;
\underbrace{\frac{2\bigl(f(w_{0})-f^{\inf}\bigr)}{\eta_{\max}T}}_{\text{optimization term}}
\;+\;
\underbrace{L\,\eta_{\max}\,\sigma^{2}}_{\text{variance term}}.
\tag{6}
\]
Consequently,
\[
\EE\!\bigl[\|\nabla f(\bar w_{T})\|^{2}\bigr]\le
\frac{2\bigl(f(w_{0})-f^{\inf}\bigr)}{\eta_{\max}T}+L\eta_{\max}\sigma^{2}.
\]
\end{theorem}

\section*{Theoretical Properties}
\begin{enumerate}[label=\arabic*.]
\item \textbf{Stability.} The bound \eqref{6} requires $\eta_{\max}\le 1/L$, which is exactly the condition guaranteeing that each SGDR step is a non‑expansive map in expectation (see Lemma~\ref{lem:descent}).

\item \textbf{Hyper‑parameter Sensitivity.}
   \begin{itemize}
   \item Larger $\eta_{\max}$ speeds up the \emph{optimization term} ($\propto 1/(\eta_{\max}T)$) but inflates the variance term $L\eta_{\max}\sigma^{2}$.
   \item The cycle length $M$ does not appear in \eqref{6} because the bound only uses the uniform upper bound $\eta_{\max}$. Hence any $M$ that respects \eqref{5} yields the same asymptotic rate.
   \end{itemize}
\item \textbf{Comparison to Classical SGD.} For a constant step size $\eta$ the standard non‑convex SGD bound is
\[
\frac{1}{T}\sum_{t=0}^{T-1}\EE\!\bigl[\|\nabla f(w_{t})\|^{2}\bigr]
\le\frac{2\bigl(f(w_{0})-f^{\inf}\bigr)}{\eta T}+L\eta\sigma^{2},
\]
which is identical to \eqref{6} with $\eta=\eta_{\max}$. Thus SGDR inherits the same worst‑case guarantee while empirically often attains faster convergence due to periodic restarts.

\item \textbf{Special Cases.}
   \begin{itemize}
   \item \emph{Deterministic case} ($\sigma=0$): the variance term vanishes and we obtain a $O(1/T)$ decay of the squared gradient norm.
   \item \emph{Exact restarts} ($M\to\infty$): the schedule reduces to simple cosine annealing without resets; the bound remains valid.
   \end{itemize}
\end{enumerate}

\section*{Discussion}
The theorem shows that SGDR achieves the same order of convergence as standard SGD for smooth non‑convex problems. The warm‑restart schedule does not deteriorate the theoretical guarantee because the analysis only requires a uniform upper bound on the step size. In practice, the periodic increase of the learning rate enables the iterates to escape shallow local minima, which explains the empirical superiority of SGDR despite identical worst‑case rates.

\newpage
\section*{Preliminaries (Definitions and Lemmas)}
\begin{lemma}[One‑step Descent under Smoothness]\label{lem:descent}
Let $w_{t+1}=w_{t}-\eta_{t}g_{t}$ with $\eta_{t}\le 1/L$. Then, conditioning on $w_{t}$,
\[
\EE\!\bigl[f(w_{t+1})\mid w_{t}\bigr]
\le f(w_{t})-\frac{\eta_{t}}{2}\,\|\nabla f(w_{t})\|^{2}
+\frac{L\eta_{t}^{2}}{2}\,\sigma^{2}.
\tag{7}
\]
\end{lemma}
\begin{proof}
\textbf{[Step 1]} Apply smoothness \eqref{3} with $y=w_{t+1}=w_{t}-\eta_{t}g_{t}$:
\[
f(w_{t+1})\le f(w_{t})-\eta_{t}\langle\nabla f(w_{t}),g_{t}\rangle
+\frac{L\eta_{t}^{2}}{2}\|g_{t}\|^{2}.
\tag{8}
\]

\textbf{[Step 2]} Take expectation w.r.t.\ $g_{t}$ conditioned on $w_{t}$.
Using \ref{as:unbiased},
\[
\EE\!\bigl[\langle\nabla f(w_{t}),g_{t}\rangle\mid w_{t}\bigr]
= \|\nabla f(w_{t})\|^{2}.
\tag{9}
\]

\textbf{[Step 3]} Decompose $\|g_{t}\|^{2}= \|\nabla f(w_{t})\|^{2}+\|g_{t}-\nabla f(w_{t})\|^{2}
+2\langle\nabla f(w_{t}),g_{t}-\nabla f(w_{t})\rangle$.
The cross term vanishes after taking conditional expectation, and by \ref{as:variance} we obtain
\[
\EE\!\bigl[\|g_{t}\|^{2}\mid w_{t}\bigr]
\le \|\nabla f(w_{t})\|^{2}+\sigma^{2}.
\tag{10}
\]

\textbf{[Step 4]} Substitute \eqref{9} and \eqref{10} into \eqref{8}:
\[
\EE\!\bigl[f(w_{t+1})\mid w_{t}\bigr]
\le f(w_{t})-\eta_{t}\|\nabla f(w_{t})\|^{2}
+\frac{L\eta_{t}^{2}}{2}\bigl(\|\nabla f(w_{t})\|^{2}+\sigma^{2}\bigr).
\tag{11}
\]

\textbf{[Step 5]} Because $\eta_{t}\le 1/L$, we have $L\eta_{t}\le 1$ and therefore
\[
-\eta_{t}+\frac{L\eta_{t}^{2}}{2}
= -\frac{\eta_{t}}{2}\bigl(2-L\eta_{t}\bigr)
\le -\frac{\eta_{t}}{2}.
\tag{12}
\]

\textbf{[Step 6]} Apply \eqref{12} to the first term of \eqref{11} and drop the non‑negative $\frac{L\eta_{t}^{2}}{2}\|\nabla f(w_{t})\|^{2}$ contribution:
\[
\EE\!\bigl[f(w_{t+1})\mid w_{t}\bigr]
\le f(w_{t})-\frac{\eta_{t}}{2}\|\nabla f(w_{t})\|^{2}
+\frac{L\eta_{t}^{2}}{2}\sigma^{2}.
\tag{13}
\]
Taking full expectation yields \eqref{7}. ∎
\end{proof}

\section*{Main Proof (Step‑by‑Step Derivation)}
We now prove Theorem~\ref{thm:main}.

\textbf{[Step 7]} Start from Lemma~\ref{lem:descent} and take total expectation:
\[
\EE\!\bigl[f(w_{t+1})\bigr]
\le \EE\!\bigl[f(w_{t})\bigr]
-\frac{\eta_{t}}{2}\EE\!\bigl[\|\nabla f(w_{t})\|^{2}\bigr]
+\frac{L\eta_{t}^{2}}{2}\sigma^{2}.
\tag{14}
\]

\textbf{[Step 8]} Rearrange \eqref{14} to isolate the gradient term:
\[
\frac{\eta_{t}}{2}\EE\!\bigl[\|\nabla f(w_{t})\|^{2}\bigr]
\le \EE\!\bigl[f(w_{t})-f(w_{t+1})\bigr]
+\frac{L\eta_{t}^{2}}{2}\sigma^{2}.
\tag{15}
\]

\textbf{[Step 9]} Sum \eqref{15} over $t=0,\dots,T-1$:
\[
\sum_{t=0}^{T-1}\frac{\eta_{t}}{2}\EE\!\bigl[\|\nabla f(w_{t})\|^{2}\bigr]
\le \EE\!\bigl[f(w_{0})-f(w_{T})\bigr]
+\frac{L\sigma^{2}}{2}\sum_{t=0}^{T-1}\eta_{t}^{2}.
\tag{16}
\]

\textbf{[Step 10]} Since $f$ is bounded below by $f^{\inf}$,
\[
\EE\!\bigl[f(w_{0})-f(w_{T})\bigr]\le f(w_{0})-f^{\inf}.
\tag{17}
\]

\textbf{[Step 11]} By the definition of the cosine schedule \eqref{1},
$0\le\eta_{t}\le\eta_{\max}$ for all $t$, thus
\[
\eta_{t}\ge\eta_{\max}\cdot\underbrace{\frac12\bigl(1+\cos(\pi)\bigr)}_{=0}
\quad\text{and}\quad
\eta_{t}^{2}\le\eta_{\max}^{2}.
\tag{18}
\]
The lower bound is zero, but for the purpose of an upper bound we keep only the uniform upper bound:
\[
\sum_{t=0}^{T-1}\eta_{t}^{2}\le T\eta_{\max}^{2}.
\tag{19}
\]

\textbf{[Step 12]} Moreover,
\[
\sum_{t=0}^{T-1}\eta_{t}\ge T\eta_{\min}\ge 0,
\]
but we need a lower bound on the *average* coefficient in front of the gradient term. Since each $\eta_{t}\le\eta_{\max}$,
\[
\frac{1}{T}\sum_{t=0}^{T-1}\eta_{t}\le\eta_{\max}.
\tag{20}
\]
To obtain a clean expression we replace each $\eta_{t}$ in the left‑hand side of \eqref{16} by the maximal admissible step size $\eta_{\max}$, which yields a *larger* left‑hand side and therefore a valid inequality:
\[
\sum_{t=0}^{T-1}\frac{\eta_{\max}}{2}\EE\!\bigl[\|\nabla f(w_{t})\|^{2}\bigr]
\le f(w_{0})-f^{\inf}
+\frac{L\sigma^{2}}{2}T\eta_{\max}^{2}.
\tag{21}
\]

\textbf{[Step 13]} Divide \eqref{21} by $T\eta_{\max}/2$:
\[
\frac{1}{T}\sum_{t=0}^{T-1}\EE\!\bigl[\|\nabla f(w_{t})\|^{2}\bigr]
\le \frac{2\bigl(f(w_{0})-f^{\inf}\bigr)}{\eta_{\max}T}
+L\eta_{\max}\sigma^{2}.
\tag{22}
\]
This is exactly the bound \eqref{6} claimed in Theorem~\ref{thm:main}. ∎

\section*{Rate Derivation (With Constants)}
From \eqref{22} we identify two contributions:
\begin{itemize}
\item \textbf{Optimization term} $O\!\bigl(1/(\eta_{\max}T)\bigr)$, which decays linearly with the number of iterations.
\item \textbf{Stochastic variance term} $L\eta_{\max}\sigma^{2}$, which is independent of $T$.
\end{itemize}
Choosing $\eta_{\max}= \Theta\!\bigl(1/\sqrt{T}\bigr)$ balances the two terms and yields the classic $O(1/\sqrt{T})$ rate for the expected gradient norm. If the variance vanishes ($\sigma=0$) the bound simplifies to $O(1/T)$.

\section*{Special Cases}
\begin{enumerate}[label=\alph*.]
\item \textbf{Deterministic Gradient Descent (no noise).} Setting $\sigma=0$ in \eqref{22} gives
\[
\frac{1}{T}\sum_{t=0}^{T-1}\|\nabla f(w_{t})\|^{2}
\le \frac{2\bigl(f(w_{0})-f^{\inf}\bigr)}{\eta_{\max}T},
\]
i.e. a $O(1/T)$ decay of the squared gradient norm.

\item \textbf{Constant Step Size without Restarts.} If we fix $\eta_{t}\equiv\eta$ and take $\eta\le 1/L$, the same derivation holds with $\eta_{\max}=\eta$, confirming that SGDR does not worsen the worst‑case guarantee.

\item \textbf{Infinite Cycle Length ($M\to\infty$).} The schedule becomes a pure cosine annealing from $\eta_{\max}$ down to $0$. Since $\eta_{t}\le\eta_{\max}$ for all $t$, the bound \eqref{22} remains valid.

\item \textbf{Finite‑variance but large variance.} When $\sigma^{2}$ dominates, the second term in \eqref{22} governs the asymptotic behaviour; the algorithm reaches a stationary neighbourhood of radius $\sqrt{L\eta_{\max}\sigma^{2}}$ around critical points.
\end{enumerate}

\end{document}